\RevisionAddBegin
The present measurements define the operating envelope. The tested settings
are interpreted as detectability-characterized annotation operating points
rather than covert-security thresholds. At 256 and 512 bits, logistic AUC
values of 0.900 and 0.957 place the current ABM configuration in a
capacity-oriented reversible-annotation regime, despite its positive net
payload and exact reversibility. At 64 bits, AUC 0.706 reflects lower
detectability, although metadata overhead dominates useful capacity. PPOCP remains the
\RevisionAddEnd
preferred evaluated method for security-prioritized operation, whereas ABM
addresses applications in which useful reversible capacity is primary and
the stego channel itself is access-controlled.

\RevisionAddBegin
The RF uniform-block experiment points to a second optimization path:
\RevisionAddEnd
coordinate coding can use enumerative subset coding conditioned on the
deterministic block order. Its grouped classification metrics show that
location selection is reliable; compressing the selected subset targets the
measured source of overhead.

\subsection{Experimental Scope}
\label{subsec:scope}

\RevisionAddBegin
The eight document images provide continuity with the binary block-mapping
literature, the disjoint BOSSbase protocol adds 100 external test images with a
four-method common evaluation, and FUNSD adds a separate frozen external test
on all 199 real scanned documents under two binarization rules. The latter
removes the earlier absence of a sizeable real scanned-document benchmark as a
primary limitation of the evidence: exact recovery and the A2-over-A1
serialization benefit are now observed without retraining on that corpus.

\paragraph{Conditions of validity.}
The exact-reversibility guarantee is conditional on a matched stego image and
its valid auxiliary stream being preserved through a lossless storage or
transmission path. The encoder may reject a requested payload when the binary
cover does not offer sufficient active mapping capacity; such rejection is an
availability outcome, not a recovery failure. A stress test on table1-3 at 256
bits confirms exact recovery after PNG and TIFF round trips, while JPEG
recompression at quality 5 and a one-pixel perturbation produce mismatched
stego/auxiliary pairs (Table~\ref{tab:channel-stress}). Applications requiring
lossy, noisy, or geometrically transformed channels therefore require an
external synchronization or error-control layer before the present exact
recovery contract applies.

\paragraph{Operating characteristics.}
Auxiliary overhead is amortized differently across image distributions and
payload sizes. On BOSSbase matched cohorts, mean net payload becomes positive
between 128 and 256 bits, whereas on FUNSD the fixed and mapping-dependent
serialization cost dominates through the tested 512-bit target. FUNSD
availability remains 98.99--100\% depending on binarization, activation, and
payload, and every available run is exact. DRD and detectability also increase
with payload. These are measured properties of the operating envelope rather
than failures of the reversible contract.

\paragraph{Remaining limitations.}
Three limitations remain material. First, the high logistic AUC at larger
BOSSbase payloads limits claims for covert or security-prioritized use under the
tested detector. Second, cross-method net-payload rankings remain partly
protocol-dependent because the baseline publications do not completely specify
standalone serialized synchronization formats. Third, although FUNSD provides
substantially stronger real-document evidence, it is one external scanned-form
corpus; generalization across additional independent document collections,
scanner pipelines, languages, and acquisition conditions remains unverified.
The present evidence also does not establish positive-net efficiency for real
scanned FUNSD documents at payloads up to 512 bits.

\RevisionAddEnd

\begin{table}[t]
\centering
\caption{Channel stress test for ABM with the same auxiliary stream.}
\label{tab:channel-stress}
\begin{tabular}{@{}lccc@{}}
\toprule
\textbf{Channel} & \textbf{Stego} &
\textbf{Message} & \textbf{Cover} \\
\midrule
PNG & Yes & Yes & Yes \\
TIFF & Yes & Yes & Yes \\
JPEG q=5 & Altered & Mismatch & Mismatch \\
One-pixel noise & Altered & Mismatch & Mismatch \\
\bottomrule
\end{tabular}
\end{table}

\section{Conclusion}
\label{sec:conclusion}

\RevisionAddBegin
This work develops and evaluates serialization-aware activation for
net-capacity-driven PEAK/ZERO block mapping in binary-image RDH. The
enumerative codec makes the active mapping description executable and charged;
the CNN ranker improves candidate placement on held-out images as a secondary
cost-ordering component. The complete system delivers 2245 net bits on average
on eight document images and 292.5 net bits on 100 BOSSbase images. A frozen
external validation on all 199 FUNSD scanned documents further shows that every
available run restores both message and binary cover exactly without training
or tuning on FUNSD.

On FUNSD, serialization-aware A2 activation reduces auxiliary cost by about
239--241 bits relative to A1 on paired available images and is never worse in
net payload under either fixed-128 or Otsu binarization. This is strong evidence
that exact reversibility and the activation benefit transfer to real scanned
documents. It is not evidence of positive net efficiency in the tested FUNSD
payload range: mean net payload remains negative through 512 bits and the
largest positive-net fraction is 2.02\%.

Under the BOSSbase common protocol, ABM obtains mean net payloads of $-57$,
$-7$, $93$, and $309$ bits at targets 64, 128, 256, and 512 bits on matched
images. Its paired net advantage is significant against every baseline after
Holm correction; at 256 bits, all adjusted $p$-values are below
$6.7\times10^{-13}$. Together, availability, binarization, runtime, memory,
and two-classifier steganalysis make this result a reproducible practical
profile beyond a single capacity number.

Under the stated implementations and wire-level accounting, the combined
evidence identifies ABM as the highest useful-capacity curve in the evaluated
common protocol, while PPOCP retains the lowest measured distortion and most
favorable detectability profile. The validity of exact recovery is conditional
on matched stego/auxiliary data over a lossless path; metadata-dominated
low-payload behavior is an operating characteristic; and the remaining true
limitations concern high-payload detectability, protocol dependence of baseline
serialization, and external breadth beyond the single FUNSD document corpus.

\RevisionAddEnd

\section*{Acknowledgements}
\RevisionAddBegin
This work was supported by the authors' institutions.
\RevisionAddEnd

\printcredits

\section*{Declarations}
\noindent\textbf{Ethics approval.}
This study does not involve human participants, animals, personal data, or
clinical interventions. It uses publicly available image data and offline
computational experiments.

\medskip
\noindent\textbf{Competing interests.}
The authors declare that they have no known competing financial interests or
\RevisionAddBegin
personal relationships that could have appeared to influence the work reported
in this paper.
\RevisionAddEnd

\medskip
\noindent\textbf{Funding.}
This research received no external funding.

\medskip
\noindent\textbf{Data availability.}
\RevisionAddBegin
BOSSbase is publicly available from its original distributor at
\url{https://dde.binghamton.edu/download/ImageDB/BOSSbase_1.01.zip} under its
applicable terms. FUNSD is publicly available from the dataset provider at
\url{https://guillaumejaume.github.io/FUNSD/dataset.zip}. The eight document
images used for block-mapping validation, derived aggregate results, per-image
evaluation artefacts, and the fixed-128/Otsu FUNSD external-validation
summaries and paired/per-image result tables accompany the reproducibility
package.
\RevisionAddEnd

\medskip
\noindent\textbf{Code availability.}
The reproducibility package contains the ABM implementation, independently
implemented comparison methods, evaluation scripts, tests, trained model
weights, CSV/JSON result tables, plotting code, and exact random seeds. It is
publicly available at
\url{https://github.com/EkodeckStephane/net-aware-block-mapping-rdh}.

